\documentclass[10.5pt, aspectratio=169]{beamer}

\usetheme{metropolis}
\usecolortheme{default}

% ここに自作スタイルファイルを読み込む
\usepackage{tikz}
\usetikzlibrary{shapes,arrows,positioning}

\usepackage{luatexja-fontspec}
\setmainjfont{Harano Aji Mincho}
\setmainjfont{Harano Aji Gothic}

% スライドのメタデータ
% ここはユーザーからの指示に応じて変更する
\title{第2章 動機づけ面接のスピリット}
\subtitle{第1部 動機づけ面接とは何か？}
\author{森口志龍}
% \institute{所属}
\date{2026-11-11}

\begin{document}

\maketitle

\begin{frame}{目次}
  \tableofcontents
\end{frame}

% ==========================================
\section{導入：なぜ「スピリット」なのか}
% ==========================================

\begin{frame}{技法（歌詞）とスピリット（曲）}
  \begin{block}{1980年代の気づき (Rollnick \& Miller, 1995)}
    \begin{itemize}
      \item 当初は「どう行うか」という\textbf{技法自体}に焦点を当てがちだった。
      \item しかし研修生の実践を見ると、「\textbf{歌詞は教えたのに曲は教えていない}」かのようだった。
      \item 根底にあるスピリット（態度と心構え）が不可欠であることに気づいた。
    \end{itemize}
  \end{block}

  \begin{alertblock}{スピリットが欠けた場合の危険性}
    \begin{itemize}
      \item 人を操作し、やりたくないことをさせる「皮相なトリック」に堕してしまう。
      \item 「間違い指摘反射」の別バージョンや、知恵で敵に勝つ知的ゲームになってしまう。
    \end{itemize}
  \end{alertblock}
\end{frame}

\begin{frame}{MIのスピリットの4つの中心要素}
  MIのスピリットは以下の4つの要素が交わるところに体現される（相互に関連し、経験的・行動的構成要素を持つ）。

  \begin{columns}[T,onlytextwidth]
    \begin{column}{0.48\textwidth}
      \begin{exampleblock}{1. パートナーシップ}
        人の「ために」、人と「共に」行う能動的な協働。
      \end{exampleblock}
      \begin{exampleblock}{2. 受容}
        相手が提示するものを心底まで受け入れ、尊重する態度。
      \end{exampleblock}
    \end{column}
    \begin{column}{0.48\textwidth}
      \begin{exampleblock}{3. 思いやり (Compassion)}
        相手の福祉を積極的に増進し、ニーズを最優先すること。
      \end{exampleblock}
      \begin{exampleblock}{4. 引き出す (Evocation)}
        相手の内側にある強み、動機、リソースを汲み出すこと。
      \end{exampleblock}
    \end{column}
  \end{columns}

  \vspace{0.5em}
  \footnotesize{※スピリットは実践前の「必須の前提条件」ではなく、\textbf{MIを実践すること自体がこれらを習慣にしてくれる}。}
\end{frame}

% ==========================================
\section{1. パートナーシップ}
% ==========================================

\begin{frame}{パートナーシップ：専門家同士の協働}
  \begin{block}{「対して」ではなく「共に」}
    \begin{itemize}
      \item 専門家が受け身の相手（まな板の上の鯉）に行うものではない。
      \item \textbf{クライエントはその人自身の専門家}である（本人以上にその人を知る人はいない）。
      \item 援助者は「同伴者」であり、忠告よりも探求、説得よりも興味と支援を行う。
    \end{itemize}
  \end{block}

  \begin{columns}[T,onlytextwidth]
    \begin{column}{0.48\textwidth}
      \begin{exampleblock}{比喩：レスリングではなく社交ダンス}
        \begin{itemize}
          \item 相手に抗うのではなく、共に動く。
          \item カウンセラーが巧みにリードを取りつつも、パートナーシップなしにダンスは成立しない。
        \end{itemize}
      \end{exampleblock}
    \end{column}
    \begin{column}{0.48\textwidth}
      \begin{alertblock}{避けるべき「専門家の罠」}
        \begin{itemize}
          \item 「カウンセラーが正しい答えをすべて持っており、それを提供するのが仕事だ」という想定を捨てる。
          \item 間違い指摘反射を意志力で抑え込む。
        \end{itemize}
      \end{alertblock}
    \end{column}
  \end{columns}
\end{frame}

\begin{frame}{パートナーシップにおける誠実さと敬意}
  \begin{itemize}
    \item \textbf{2つの願望の調和}
      \begin{itemize}
        \item クライエントの願望だけでなく、カウンセラー自身の願望（アジェンダ・価値観）についても意識し、\textbf{正直であること}が必須。
      \end{itemize}
    \vspace{0.5em}
    \item \textbf{深い尊敬と特権}
      \begin{itemize}
        \item MIの会話は「ソファーで隣に座り、相手の誕生から今までの写真アルバムを一緒に見ること」に似ている。
        \item カウンセラーの目的は、自分の世界観を披露することではなく、\textbf{相手の目線から世界を見ること}。
      \end{itemize}
  \end{itemize}
\end{frame}

% ==========================================
\section{2. 受容}
% ==========================================

\begin{frame}{受容（Acceptance）の4つの側面}
  「受容」とは行為の承認や現状の黙認ではない。Carl Rogersの仕事に源流を持つ4つの側面から構成される。

  \begin{columns}[T,onlytextwidth]
    \begin{column}{0.48\textwidth}
      \begin{block}{① 絶対的価値 (Absolute Worth)}
        \begin{itemize}
          \item 無条件の肯定的関心・非支配的な配慮。
          \item \textbf{反対：}条件付きの価値づけ（判断の態度）。
          \item ありのままで受容されると自由な変化が可能になる（自己実現傾向・テロス）。
        \end{itemize}
      \end{block}
    \end{column}
    \begin{column}{0.48\textwidth}
      \begin{block}{② 正確な共感 (Accurate Empathy)}
        \begin{itemize}
          \item 相手の内的世界を「あたかも」自分自身のものであるかのように理解する努力（同情や同一化ではない）。
          \item \textbf{反対：}自分の視点の押しつけ。
        \end{itemize}
      \end{block}
    \end{column}
  \end{columns}
\end{frame}

\begin{frame}
  \begin{columns}[T, onlytextwidth]
    \begin{column}{0.48\textwidth}
      \begin{block}{③ 自律性のサポート (Autonomy)}
        \begin{itemize}
          \item 自身の方向を決める奪えない権利と能力の尊重。
          \item \textbf{反対：}強制とコントロール（心理的反動・防衛性を生む）。
          \item 選択の自由を認めると防衛性が減り、変化が促進される。
        \end{itemize}
      \end{block}
    \end{column}
    \begin{column}{0.48\textwidth}
      \begin{block}{④ 是認 (Affirmation)}
        \begin{itemize}
          \item 相手の強みと努力を探し出して承認する意図的なコミュニケーション。
          \item \textbf{反対：}悪い点（欠陥）を探し、治し方を教えること。
        \end{itemize}
      \end{block}
    \end{column}  
  \end{columns}
\end{frame}

% ==========================================
\section{3. 思いやり}
% ==========================================

\begin{frame}{思いやり（Compassion）}
  第3版から新たに加えられた、MIの基盤となる第3の要素。

  \begin{block}{思いやりの定義}
    \begin{itemize}
      \item 単なる感情的体験（共に苦しむこと、同情）ではない。
      \item \textbf{相手の福祉を積極的に増進しようとすることであり、相手のニーズを満たすことを優先する意図的なコミットメント。}
    \end{itemize}
  \end{block}

  \begin{alertblock}{なぜ「思いやり」が追加されたのか？}
    \begin{itemize}
      \item 他の3要素（パートナーシップ・受容・引き出す）は、\textbf{自分の私的利益を得る目的（有能な販売員や搾取など）でも実践可能}だからである。
      \item 治療者が育む信頼が適切なものになるよう、クライエントの最善の利益を第一とする心構えが不可欠。
    \end{itemize}
  \end{alertblock}
\end{frame}

% ==========================================
\section{4. 引き出す}
% ==========================================

\begin{frame}{引き出す（Evocation）：欠陥モデルからの脱却}
  \begin{columns}[T,onlytextwidth]
    \begin{column}{0.48\textwidth}
      \begin{alertblock}{従来の「欠陥モデル」 (\textit{docere})}
        \begin{itemize}
          \item 前提：「あなたに欠けているものを私が与えよう」
          \item 知識、洞察、対処技能などを外から押し込む（講義型）。
          \item 感染症治療や車の修理には有効だが、\textbf{個人の変化には機能しにくい}。
        \end{itemize}
      \end{alertblock}
    \end{column}
    \begin{column}{0.48\textwidth}
      \begin{exampleblock}{MIの「引き出すモデル」 (\textit{e ducere})}
        \begin{itemize}
          \item 前提：「必要なものはあなたが持っています。一緒に発見しましょう」
          \item 井戸から水を汲み上げるように、内なる知恵と経験を汲み出す。
          \item 強みとリソースにフォーカスする。
        \end{itemize}
      \end{exampleblock}
    \end{column}
  \end{columns}

  \vspace{0.8em}
  \begin{block}{両価性（Ambivalence）との合致}
    変化に両価的な人の心（内部委員会）には、\textbf{すでに変化に賛成の意見が存在している}。治療者の仕事は、最初から存在する変化への動機づけを誘い出し、強化することである。
  \end{block}
\end{frame}

% ==========================================
\section{まとめ：人間中心ケアの原則と成長}
% ==========================================

\begin{frame}{MIの実践的定義とパーソン・センタード・ケアの原則}
  \begin{exampleblock}{臨床家のためのMIの定義（第2の定義）}
    「動機づけ面接は\textbf{パーソン・センタード・カウンセリング・スタイル}であり、変化に関する\textbf{両価性}から生じる一般的な問題を扱うものである。」
  \end{exampleblock}

  \textbf{パーソン・センタード・ケア（人間中心ケア）の主な原則：}
  \begin{enumerate}
    \item サービスは受益者のために存在し、クライエントのニーズが優先される。
    \item 変化とは本質的に\textbf{自己変化}であり、クライエントは自分自身の専門家である。
    \item 変化を起こす責任は臨床家にはない（臨床家だけでは何もできない）。
    \item 会話は権力闘争（レスリング）ではなく、\textbf{ダンス}のように感じられるべきである。
    \item 動機づけは外から植えつけるものではなく、\textbf{中から引き出されるもの}である。
    \item 行動に関する選択の最終決断は、常にクライエント自身にある。
  \end{enumerate}
\end{frame}

\begin{frame}{成長のプロセスとキーポイント}
  \begin{block}{実践を通じた成長}
    \begin{itemize}
      \item スピリットを最初から完全に体得している必要はない。
      \item 知識を振り返り、他者といるときにこのスタイルを実践し続けることで、心構えが内面化され、自発的・自動的なものになっていく。
    \end{itemize}
  \end{block}

  \begin{exampleblock}{本章のキーポイント}
    \begin{itemize}
      \item MIは誰かに「対して」ではなく、誰かの「\textbf{ために}」、誰かと「\textbf{共に}」行われる。
      \item MIのスピリットの4特質：\textbf{パートナーシップ、受容、思いやり、引き出す}。
      \item 受容の4側面：\textbf{絶対的価値、正確な共感、自律性のサポート、是認}。
      \item MIはもともと存在しているものを引き出すことであり、欠けているものを植えつけることではない。
    \end{itemize}
  \end{exampleblock}
\end{frame}

\end{document}